\begin{abstract}

\heng{need to rewrite so it's direct and clear what prolems you are studying: how incorrect information travel inside LLM? and summarize your findings}
Post-training adapts large language models using massive, heterogeneous data mixtures, where some noise is unavoidable and often too subtle to reliably filter. \heng{previous sentence is terse and difficult to parse. rewrite it} Diagnosing its impact is difficult: practitioners rarely isolate causal examples, \heng{not sure what you mean by 'isolate causal examples'} and internally an LLM’s knowledge is highly entangled, so a small corruption may have non-local effects. We study two consequences of noisy post-training knowledge \heng{what is noisy post-training knowledge}: \emph{knowledge confidence} (whether model confidence remains a useful proxy for correctness) and the \emph{knowledge ripple effect} (how corruption spreads to related facts). We propose a controlled framework that injects structured false knowledge during post-training and probes accuracy and token-level confidence on the injected facts and their graph neighbors. \heng{you never talked about how the graph is constructed} Surprisingly, small amounts of structured noise \heng{what is structured noise?} can make models confidently wrong and induce measurable spillover to nearby knowledge, with topology shaping the blast radius. These findings caution against confidence-only reliability signals and motivate structure-aware evaluation for post-training and knowledge editing.
\end{abstract}